\documentclass[10pt,a4paper]{amsart}
\usepackage[margin=19mm]{geometry}
\usepackage{amsmath,amssymb,mathtools}
\usepackage[scr=rsfs,cal=euler]{mathalfa}
\usepackage{xfrac}
\usepackage[unicode,hidelinks,bookmarks=false]{hyperref}
\newcommand{\ie}{i.e.\ }
\newcommand{\cf}{cf.\ }
\newcommand{\resp}{respectively\ }
\newcommand{\Subsection}{\subsection}
\providecommand{\nobreakdash}{\nobreak\hbox{-}}
\setlength{\emergencystretch}{2em}
\multlinegap0pt
\usepackage{bm}
\usepackage{bbm}
\usepackage{dsfont}
\usepackage{xspace}
\usepackage{cancel}
\usepackage{mathtools}
\usepackage{amsthm}
\makeatletter
\@ifundefined{theorem}{\newtheorem{theorem}{Theorem}}{}
\renewenvironment{proof}[1]{\par\medskip\noindent\textbf{#1}\enspace\ignorespaces}{\par\medskip}
\@ifundefined{lemma}{\newtheorem{lemma}[theorem]{Lemma}}{}
\@ifundefined{prop}{\newtheorem{prop}[theorem]{Proposition}}{}
\makeatother
\title[On projective threefolds with two-dimensional space of vanis]{On projective threefolds with two-dimensional space of vanishing cycles}
\author{Timofey Fedorov}
\date{Source: arXiv:2501.13027v3 (CC BY 4.0)}
\begin{document}
\maketitle

\noindent\emph{Source and licence.} On projective threefolds with two-dimensional space of vanishing cycles. Version arXiv:2501.13027v3 (CC BY 4.0), distributed under the Creative Commons Attribution 4.0 International licence, as recorded in the arXiv licence metadata for this exact version and shown on the abstract page https://arxiv.org/abs/2501.13027v3. Full rights notice: licenses/arxiv-2501.13027-CC-BY-4.0.txt.\par\smallskip

\noindent\textit{Editorial scope.} Counted unit: the Lemma whose source label is \texttt{obvious} in the article (source0.tex lines 267--271, source label \texttt{obvious}), delivered together with the article's own complete original proof of it (source0.tex lines 273--422, one \texttt{proof} environment of 150 lines). No mathematics is written, repaired or re-derived: statement and proof are the article's own text line for line. Required context retained uncounted: 2-jet (lines 261--265). Every definition, display and label of those lines is retained unchanged. Omitted from this package and not redistributed: 1--260, 266--266, 272--272, 423--632. Nothing inside a retained excerpt is omitted or altered: every retained body is delivered exactly as the article's lines are written, so no omission receipt of any kind accompanies this package. Citations: the retained excerpts carry no citation command at all, so no bibliography entry of the article is transcribed and no reference is redistributed. Reference adaptation: none was needed and none was made; every reference inside the delivered unit resolves inside it, so no unresolved reference or citation is produced. Printed slips kept literal: no printed word, symbol, hypothesis or number of the counted statement or of its proof was altered. Layout and class: the article's own class is \texttt{amsart} with packages including cmap, fontenc, inputenc, amsfonts, amsmath, mathtools, upgreek, graphicx, amssymb, amsthm, enumitem, esint, hyperref, comment; the wrapper is the lane's pinned \texttt{amsart} editorial wrapper at 10pt on A4, which restates the theorem-like environments the retained text needs and adds the article's own macro definitions that the retained text needs (none), with extra wrapper packages (bm, bbm, dsfont, xspace, cancel, mathtools). The delivered numbering is the unit's own. Assets: no retained span contains a figure environment, a graphics-inclusion command, a PDF-page inclusion, a table or a TikZ picture; no image file is copied into this package, the package declares no asset dependency and no image file is redistributed with it. Licence: version arXiv:2501.13027v3 of this article is distributed under the Creative Commons Attribution 4.0 International licence, as recorded in the arXiv licence metadata for this exact version and shown on the abstract page https://arxiv.org/abs/2501.13027v3. A full rights notice is delivered at licenses/arxiv-2501.13027-CC-BY-4.0.txt. Characters: every retained excerpt is pure ASCII, so the delivered engine, XeLaTeX with its default Latin Modern text font, reports no missing character.
\begin{theorem}
    \label{2-jet}
    Suppose that $E$ is a very ample rank-$2$ bundle on a surface $S$.
    Then $\mathrm{det}\, E$ is $2$-jet ample.
\end{theorem}

\begin{lemma}
    \label{obvious}
    Let $S$ and $E$ be as before. Take pairwise disctinct points $y, y_1,\dots, y_k$ and an arbitrary point $p \in s_{y_1}$.
    Then there is a hyperplane $H \subset \mathbb P^n$ containing $s_y$, not containing $s_{y_2}, \dots, s_{y_k}$ (hence, intersecting each of these fibers at one point) and intersecting the fiber $s_{y_1}$ at the point $p$.\qed
\end{lemma}

\begin{proof}{(of Theorem \ref{2-jet})}
    %We will only need sections from $H^0(\mathrm{det}\, E)$ that are linear combinations of elements of the form $g_1 \wedge g_2, g_i \in H^0(E)$ to separate $2$-jets on $S$.

    Consider $r\le3$ distinct points $y_1,\dots,y_r\in S$ and $r$ natural numbers $k_1,\dots,k_r$, $\sum k_i=3$.
    We need to examine three different cases:

    (1) $r=3$, $k_1=k_2=k_3=1$. We prove surjectivity of
    \[
    \varphi \colon H^0(\mathrm{det}\, E) \to H^0(\mathrm{det} \,E\otimes (\mathcal O_{y_1}/ \mathfrak m_{y_1}\oplus \mathcal O_{y_2}/ \mathfrak m_{y_2} \oplus \mathcal O_{y_3}/ \mathfrak m_{y_3})) \cong \Bbbk^3
    \]
    as follows. Choose a section $\alpha_1 \in H^0(E)$ such that $H_{\alpha_1} \supset s_{y_1}$ and $H_{\alpha_1}$ intersects each fiber $s_{y_2}, s_{y_3}$ at one point (using lemma \ref{obvious}); let $p$ be the intersection point of $H_{\alpha_1}$ and $s_{y_3}$.
    Choose a section $\alpha_2 \in H^0(E)$ such that $H_{\alpha_2} \supset s_{y_2}$, $H_{\alpha_2}$ intersects each fiber $s_{y_1}, s_{y_3}$ at one point and intersection of $H_{\alpha_2}$ with $s_{y_3}$ is not the point $p$.
    Then $\varphi (\alpha_1 \wedge \alpha_2)$ is equal to $(0, 0, \lambda), \lambda \neq 0$. Hence all standard basis elements are in the image of $\varphi$ and we are done.

    (2) $r=2$, $k_1=2$, $k_2=1$. We need to prove surjectivity of
    \[
    \varphi \colon H^0(\mathrm{det}\, E) \to H^0(\mathrm{det} \,E \otimes (\mathcal O_{y_1}/\mathfrak m_{y_1}^2\oplus \mathcal O_{y_2} / \mathfrak m_{y_2})) \cong \Bbbk^3\oplus \Bbbk
    \]
    Fix an affine neighborhood $U \ni y_1$ such that $E$ is trivial over $U$. Choose local coordinates $x_1, x_2$ in $\mathcal O_{S, y_1}$.
    Every $f\in H^0(U, E)$ is represented as the pair $f = (f_1, f_2), f_i \in \mathcal O_{S,y_1}$ and since $O_{S,y_1}$ is a regular local ring the images of $f_1, f_2$ in the completion $\widehat{\mathcal O_{S, y_1}} \cong \Bbbk[[x_1, x_2]]$ are represented by the power series $f_1(x_1, x_2), f_2(x_1, x_2)$ in variables $x_1, x_2$.
    Very ampleness of $\mathcal O(1)$ means, in particular, that $\mathcal O(1)$ generates $1$-jets in all points $p \in s_{y_1}$, i.e., for any $p \in s_{y_1}$ the natural map
    \[
    \psi_p \colon H^0(\mathcal O(1)) \to H^0(\mathcal O(1) \otimes \mathcal O_{X, p}/ \mathfrak m_{X, p}^2)
    \]
    is surjective.
    We want to rewrite these conditions for $\psi_p,\, p \in s_{y_1}$ in terms of sections $f\in H^0(E)$.
    Let $(t_1 \colon t_2)$ be homogeneous coordinates on $\mathbb P^1$, so $((t_1 \colon t_2), x_1, x_2)$ are coordinates for $\pi^{-1}(U) \cong \mathbb P^1 \times U$.
    Section $f\in H^0(E)$ written locally as $(f_1(x_1, x_2), f_2(x_1, x_2))$ correspond to the section $\bar f\in H^0(\mathrm{det}\, E)$ written as
    \[
    ((t_1:t_2), x_1, x_2) \mapsto t_1 f_2(x) - t_2 f_1(x),\, x = (x_1, x_2) \in U
    \]
    Consider sections $f \in H^0(E)$ such that $(f_1(y_1)\colon f_2(y_1)) = (1:\lambda), \lambda \in \Bbbk$ and take $u = t_2/t_1$ to be the affine coordinate near $(1:\lambda) \in \mathbb P^1$.
    In coordinates $(u, x_1, x_2)$ section $\bar f$ is written as
    \[
    (u, x_1, x_2) \mapsto f_2(x)-uf_1(x)
    \]
    The gradient vector of $\bar f$ w.r.t $u, x_1, x_2$ equals
    \[
    \left(-f_1(x), \frac{\partial f_2}{\partial x_1} - u\frac{\partial f_1}{\partial x_1}, \frac{\partial f_2}{\partial x_2} - u\frac{\partial f_1}{\partial x_2}\right)
    \]
    Therefore, we can restate the surjectivity condition of all $\psi_p$ as follows:
    \begin{prop}
        \label{matrices}
        Let $\psi_p, p \in s_{y_1}$ be the natural maps as above.
        All the maps $\psi_p$ are surjective iff the following holds: differentials of sections $f\in H^0(E)$ at the point $y_1$ form a linear subspace $\Pi\subset \mathfrak m_{y_1}/\mathfrak m_{y_1}^2 \oplus \mathfrak m_{y_1}/\mathfrak m_{y_1}^2 \cong \mathrm{Mat}_{2 \times 2}(\Bbbk)$
        such that for any nonzero vector $v = (t_1, t_2)^T \in \Bbbk^2$ there exist two matrices $A_1, A_2 \in \Pi$, for which vectors $A_1 v$ and $A_2 v$ are linearly independent,
        and for any $w \in \Bbbk^2$ and $A\in \Pi$ one can choose a section $f \in H^0(E)$ satisfying $f(y_1) = w, (df)_{y_1} = A$.\qed
    \end{prop}
    
    To complete the proof in the case (2) choose a section $g \in H^0(E)$ such that $g(y_1) = (1, 0)$ and the corresponding hyperplane contains fiber $s_{y_2}$ (i.e. $g(y_2) = 0$) using Lemma \ref{obvious}.
    Consider sections $f \in H^0(E)$ with $f(y_1) = 0$. Write power series representations of $g$ and $f$ in a $2$ by $2$ matrix $J$:
    \begin{equation*}
        \begin{pmatrix}
            1+g_1(x_1,x_2)&g_2(x_1,x_2)\\
            f_1(x_1,x_2)  &f_2(x_1,x_2)\\
        \end{pmatrix}
    \end{equation*}
    Here the power series $g_i, f_i$ are in $\widehat{m_{y_1}} \subset \widehat{\mathcal O_{S, y_1}}$.
    The differential of $g\wedge f$ at the point $y_1$ is the linear part of the power series $\mathrm{det}\, J$, so it is actually equal to the linear part of $f_2$.
    Due to Proposition \ref{matrices} we can choose $f$ such that section $g\wedge f$ have an arbitrary differential at $y_1$.
    Therefore the image of $\varphi$ contains the subspace $(\mathfrak m_{y_1}/\mathfrak m_{y_1}^2, 0) \subset \mathcal O_{y_1} / \mathfrak m_{y_1}^2 \oplus \mathcal O_{y_2}/ \mathfrak m_{y_2}$.
    Similarly to the case (1) one can a get vector of the form $(1+ g(x_1, x_2), 0)$ and the vector $(0, 1)$, hence $\varphi$ is surjective.
    
    (3) $r=1$, $k_1=3$; we will write $y$ instead of $y_1$. We are to prove the surjectivity of
    \[
    \varphi \colon H^0(\mathrm{det}\, E) \to H^0(\mathrm{det} \,E/ \mathfrak m_{y}^3)
    \]
    Let us analyse possible subspaces $\Pi \subset \mathrm{Mat}_{2 \times 2}(\Bbbk)$ satisfying conditions of Proposition \ref{matrices}.
    Note that if $\Pi$ is equal to $\mathrm{Mat}_{2 \times 2}(\Bbbk)$ there is nothing to prove: if one consider sections $f, g \in H^0(E)$ having zero at $y$, then quadratic part of $f\wedge g$ spans $\mathfrak m_y^2/\mathfrak m_y^3$.
    Combining this with arguments from the case (2) yields surjectivity of $\varphi$.
    
    Note also that $\Pi$ cannot be $2$-dimensional.
    Indeed, if $\Pi$ is spanned by $A_1, A_2 \in \mathrm{Mat}_{2 \times 2}(\Bbbk)$, then the hypothesis of Proposition \ref{matrices} implies that for any nonzero vector $v = (t_1, t_2)^T$ vectors $A_1 v, A_2 v$ are linearly independent.
    The determinant of the matrix constructed from two vectors $A_1 v, A_2 v$ is equal to
    \[
    t_1^2(a_1 c_2-a_2 c_1) +t_1 t_2 (a_1 d_2 - b_2 c_1 + b_1 c_2 - a_2 d_1) +t_2^2(b_1 d_2 - b_2 d_1)
    \]
    Clearly there exist a pair $(t_1, t_2)$ such that this expression vanishes, thus $\Pi$ cannot be of dimension two.

    Now consider the case of $3$-dimensional $\Pi$. We want to choose a suitable basis for $\Pi$.
    By Proposition \ref{matrices} there are two matrices $A_1, A_2$ in $\Pi$ such that $A_1 (1, 0)^T, A_2 (1, 0)^T$ are linearly independent. Hence we can choose $A_1, A_2$ of the form
    {\sloppy

    }
    \begin{equation*}
        A_1 =\begin{pmatrix}
            1&a_1\\
            0&b_1\\
        \end{pmatrix},\,
        A_2 =\begin{pmatrix}
            0&a_2\\
            1&b_2\\
        \end{pmatrix}
    \end{equation*}
    There exists a matrix $A_3 \in \Pi \setminus \mathrm{Span}\, (A_1, A_2)$ of the form
    \begin{equation*}
        A_3 =\begin{pmatrix}
            a_3&1\\
            b_3&0\\
        \end{pmatrix}
    \end{equation*}
    or
    \begin{equation*}
        A_3 =\begin{pmatrix}
            a_3&0\\
            b_3&1\\
        \end{pmatrix}
    \end{equation*}
    and after swapping the rows in all the matrices in $\Pi$ if necessary we can assume that $A_3$ has the form
    \begin{equation*}
        A_3 =\begin{pmatrix}
            a_3&1\\
            b_3&0\\
        \end{pmatrix}
    \end{equation*}
    We want to prove that sections of the form $f\wedge g \in H^0(\mathrm{det}\, E)$, $f(y) = g(y) = 0$, span $\mathfrak m_y^2/\mathfrak m_y^3$,
    that is, that the three homogeneous quadratic polynomials 
    \begin{align*}
    (A_1 \cdot (t_1, t_2)^T) &\wedge (A_2 \cdot (t_1, t_2)^T)\\
    (A_2 \cdot (t_1, t_2)^T) &\wedge (A_3 \cdot (t_1, t_2)^T)\\
    (A_1 \cdot (t_1, t_2)^T) &\wedge (A_3 \cdot (t_1, t_2)^T)\\
    \end{align*}
    in variables $t_1, t_2$ are linearly independent.
    We show this by direct computation using Maple.
    Writing these three polynomials into the $3$-by-$3$ matrix (the first polynomial in the first row and so on) with respect to the basis $t_1^2, t_1 t_2, t_2^2$ we have:
    \begin{equation*}
        D =\begin{pmatrix}
            1&a_1 + b_2 &a_1 b_2 - a_2 b_1  \\
            b_3&a_1 b_3 - b_1 a_3&-b_1      \\
            -a_3& a_2 b_3 - a_3 b_2 -1&-b_2 \\
        \end{pmatrix}
    \end{equation*}
    Note that $P_1$ does not vanish at the point $(t_1\colon t_2) = (1\colon 0)$. Putting $t_2 = 1$ we get polynomials in variable $t_1$:
    \begin{align*}
        P_1 &= t_1^2 + t_1(a_1 +b_2) + a_1 b_2 - a_2 b_1\\
        P_2 &= t_1^3 b_3 + t_1 (a_1 b_3 - b_1 a_3) - b_1\\
        P_3 &= t_1^3 (-a_3) +t_1(a_2 b_3 - a_3 b_2 -1) - b_2\\
    \end{align*}
    We have
    \begin{align*}
        \mathrm{det}\, D &= a_1^2 a_3 b_2 b_3 - a_1 a_2 a_2 b_1 b_3 + a_1 a_2 b_2 b_3^2 - a_1 a_3^2 b_1 b_2 - a_1 a_3 b_2^2 b_3 -\\
        &- a_2^2 b_1 b_3^2+a_2 a_3^2 b_1^2 + a_2 a_3 b_1 b_2 b_3 + a_1 a_3 b_1 -\\
        &-a_1 b_2 b_3 + 2 a_2 b_1 b_3 + a_3 b_1 b_2 + b_2^2 b_3 - b_1\\
    \end{align*}
    and $res(P_1, P_2) = -b_1 \cdot \mathrm{det} \, D$, $res(P_1, P_3) = a_2 \cdot \mathrm{det}\, D$. %(note that $-b_1 = \mathrm{det}\, A_1$, $a_2 = -\mathrm{det}\, A_2$, so these equations rewrites as )
    Suppose now that $\mathrm{det}\, D$ is zero. Then $res(P_1, P_2) = res(P_1, P_3) = 0$.
    Since the rows of the matrix $D$ are linearly dependent, either the second row is a linear combination of the first and the third rows or the third row is a linear combination of the first and the second row.
    In both cases all three polynomials have a common zero contradicting Proposition \ref{matrices}.
    Thus $\mathrm{det}\, D \neq 0$ and we are done.
\end{proof}

\end{document}
